\documentclass[11pt,a4paper]{article}
% Shared preamble for the qi26_25 weekly deliverable reports.
%
% Each weekly report is a short standalone document that inputs this file, so
% that formatting stays consistent across all six without duplicating the
% package list in every one. Change something here and every report follows.
%
% Usage in a weekly report:
%     \documentclass[11pt,a4paper]{article}
%     % Shared preamble for the qi26_25 weekly deliverable reports.
%
% Each weekly report is a short standalone document that inputs this file, so
% that formatting stays consistent across all six without duplicating the
% package list in every one. Change something here and every report follows.
%
% Usage in a weekly report:
%     \documentclass[11pt,a4paper]{article}
%     % Shared preamble for the qi26_25 weekly deliverable reports.
%
% Each weekly report is a short standalone document that inputs this file, so
% that formatting stays consistent across all six without duplicating the
% package list in every one. Change something here and every report follows.
%
% Usage in a weekly report:
%     \documentclass[11pt,a4paper]{article}
%     \input{weekly_preamble}
%     \weekhead{1}{Vision Foundation and Data Encoding}{Week of ...}
%     ... body ...
%     \end{document}

\usepackage[utf8]{inputenc}
\usepackage[T1]{fontenc}
\usepackage[margin=2.5cm]{geometry}
\usepackage{amsmath,amssymb}
\usepackage{booktabs}
\usepackage{graphicx}
\usepackage{enumitem}
\usepackage{listings}
\usepackage{xcolor}
\usepackage[hidelinks]{hyperref}
\usepackage{titlesec}
\usepackage{fancyhdr}

\definecolor{codebg}{RGB}{247,247,250}
\definecolor{codekw}{RGB}{20,90,160}
\definecolor{codecm}{RGB}{110,120,130}
\definecolor{accent}{RGB}{20,90,160}

\lstset{
  backgroundcolor=\color{codebg},
  basicstyle=\ttfamily\footnotesize,
  keywordstyle=\color{codekw}\bfseries,
  commentstyle=\color{codecm}\itshape,
  breaklines=true,
  frame=single,
  rulecolor=\color{gray!30},
  language=Python,
  showstringspaces=false,
  xleftmargin=6pt,
  framexleftmargin=6pt
}

\setlist{itemsep=2pt,topsep=4pt}
\titleformat{\section}{\normalfont\large\bfseries\color{accent}}{\thesection}{0.8em}{}
\titleformat{\subsection}{\normalfont\normalsize\bfseries}{\thesubsection}{0.8em}{}

\pagestyle{fancy}
\fancyhf{}
\fancyfoot[C]{\small\thepage}
\fancyfoot[R]{\scriptsize\texttt{github.com/Pushkar0997/qi26\_25}}
\renewcommand{\headrulewidth}{0pt}

% \weekhead{number}{title}{deliverable as stated in the project brief}
\newcommand{\weekhead}[3]{%
  \begin{center}
    {\scriptsize\textsf{QINTERN 2026 \textbullet\ QWORLD \textbullet\ WEEKLY DELIVERABLE}}\\[0.25cm]
    {\LARGE\bfseries Week #1: #2}\\[0.35cm]
    {\normalsize Hybrid QML and Quantum String Algorithms for Document Extraction}\\[0.25cm]
    {\small Pushkar Kumar \quad\textbullet\quad Mentor: Potluri Krishna Priyatham}\\[0.4cm]
    \fbox{\begin{minipage}{0.88\textwidth}
      \small\textbf{Deliverable as specified in the project brief:} #3
    \end{minipage}}
  \end{center}
  \vspace{0.4cm}
}

% A boxed note for caveats and honest limitations.
\newcommand{\honestnote}[1]{%
  \vspace{0.2cm}
  \noindent\fcolorbox{gray!40}{gray!8}{%
    \begin{minipage}{0.97\textwidth}\small #1\end{minipage}}
  \vspace{0.2cm}
}

%     \weekhead{1}{Vision Foundation and Data Encoding}{Week of ...}
%     ... body ...
%     \end{document}

\usepackage[utf8]{inputenc}
\usepackage[T1]{fontenc}
\usepackage[margin=2.5cm]{geometry}
\usepackage{amsmath,amssymb}
\usepackage{booktabs}
\usepackage{graphicx}
\usepackage{enumitem}
\usepackage{listings}
\usepackage{xcolor}
\usepackage[hidelinks]{hyperref}
\usepackage{titlesec}
\usepackage{fancyhdr}

\definecolor{codebg}{RGB}{247,247,250}
\definecolor{codekw}{RGB}{20,90,160}
\definecolor{codecm}{RGB}{110,120,130}
\definecolor{accent}{RGB}{20,90,160}

\lstset{
  backgroundcolor=\color{codebg},
  basicstyle=\ttfamily\footnotesize,
  keywordstyle=\color{codekw}\bfseries,
  commentstyle=\color{codecm}\itshape,
  breaklines=true,
  frame=single,
  rulecolor=\color{gray!30},
  language=Python,
  showstringspaces=false,
  xleftmargin=6pt,
  framexleftmargin=6pt
}

\setlist{itemsep=2pt,topsep=4pt}
\titleformat{\section}{\normalfont\large\bfseries\color{accent}}{\thesection}{0.8em}{}
\titleformat{\subsection}{\normalfont\normalsize\bfseries}{\thesubsection}{0.8em}{}

\pagestyle{fancy}
\fancyhf{}
\fancyfoot[C]{\small\thepage}
\fancyfoot[R]{\scriptsize\texttt{github.com/Pushkar0997/qi26\_25}}
\renewcommand{\headrulewidth}{0pt}

% \weekhead{number}{title}{deliverable as stated in the project brief}
\newcommand{\weekhead}[3]{%
  \begin{center}
    {\scriptsize\textsf{QINTERN 2026 \textbullet\ QWORLD \textbullet\ WEEKLY DELIVERABLE}}\\[0.25cm]
    {\LARGE\bfseries Week #1: #2}\\[0.35cm]
    {\normalsize Hybrid QML and Quantum String Algorithms for Document Extraction}\\[0.25cm]
    {\small Pushkar Kumar \quad\textbullet\quad Mentor: Potluri Krishna Priyatham}\\[0.4cm]
    \fbox{\begin{minipage}{0.88\textwidth}
      \small\textbf{Deliverable as specified in the project brief:} #3
    \end{minipage}}
  \end{center}
  \vspace{0.4cm}
}

% A boxed note for caveats and honest limitations.
\newcommand{\honestnote}[1]{%
  \vspace{0.2cm}
  \noindent\fcolorbox{gray!40}{gray!8}{%
    \begin{minipage}{0.97\textwidth}\small #1\end{minipage}}
  \vspace{0.2cm}
}

%     \weekhead{1}{Vision Foundation and Data Encoding}{Week of ...}
%     ... body ...
%     \end{document}

\usepackage[utf8]{inputenc}
\usepackage[T1]{fontenc}
\usepackage[margin=2.5cm]{geometry}
\usepackage{amsmath,amssymb}
\usepackage{booktabs}
\usepackage{graphicx}
\usepackage{enumitem}
\usepackage{listings}
\usepackage{xcolor}
\usepackage[hidelinks]{hyperref}
\usepackage{titlesec}
\usepackage{fancyhdr}

\definecolor{codebg}{RGB}{247,247,250}
\definecolor{codekw}{RGB}{20,90,160}
\definecolor{codecm}{RGB}{110,120,130}
\definecolor{accent}{RGB}{20,90,160}

\lstset{
  backgroundcolor=\color{codebg},
  basicstyle=\ttfamily\footnotesize,
  keywordstyle=\color{codekw}\bfseries,
  commentstyle=\color{codecm}\itshape,
  breaklines=true,
  frame=single,
  rulecolor=\color{gray!30},
  language=Python,
  showstringspaces=false,
  xleftmargin=6pt,
  framexleftmargin=6pt
}

\setlist{itemsep=2pt,topsep=4pt}
\titleformat{\section}{\normalfont\large\bfseries\color{accent}}{\thesection}{0.8em}{}
\titleformat{\subsection}{\normalfont\normalsize\bfseries}{\thesubsection}{0.8em}{}

\pagestyle{fancy}
\fancyhf{}
\fancyfoot[C]{\small\thepage}
\fancyfoot[R]{\scriptsize\texttt{github.com/Pushkar0997/qi26\_25}}
\renewcommand{\headrulewidth}{0pt}

% \weekhead{number}{title}{deliverable as stated in the project brief}
\newcommand{\weekhead}[3]{%
  \begin{center}
    {\scriptsize\textsf{QINTERN 2026 \textbullet\ QWORLD \textbullet\ WEEKLY DELIVERABLE}}\\[0.25cm]
    {\LARGE\bfseries Week #1: #2}\\[0.35cm]
    {\normalsize Hybrid QML and Quantum String Algorithms for Document Extraction}\\[0.25cm]
    {\small Pushkar Kumar \quad\textbullet\quad Mentor: Potluri Krishna Priyatham}\\[0.4cm]
    \fbox{\begin{minipage}{0.88\textwidth}
      \small\textbf{Deliverable as specified in the project brief:} #3
    \end{minipage}}
  \end{center}
  \vspace{0.4cm}
}

% A boxed note for caveats and honest limitations.
\newcommand{\honestnote}[1]{%
  \vspace{0.2cm}
  \noindent\fcolorbox{gray!40}{gray!8}{%
    \begin{minipage}{0.97\textwidth}\small #1\end{minipage}}
  \vspace{0.2cm}
}


\begin{document}
\weekhead{2}{QML for Feature Extraction (OCR Stage)}
{Simulation results of the quantum feature extraction layer.}

\section{Objective}

Build a quantum neural network component for character and layout detection:
design a quanvolutional layer to identify document features, and train a
variational quantum circuit for basic character recognition.

\section{The Quanvolutional Filter}

A quanvolutional layer is the quantum analogue of a convolution. A small window
slides over the image; at each position the patch is encoded into a quantum
state, transformed by a circuit, and measured.

The filter is a fixed 4-qubit circuit: two layers of $R_Y$ rotations with a ring
of CX gates between them, angles drawn once from a seeded generator and held
constant. This is the canonical \emph{untrained random} formulation of Henderson
et al.

Two architectural points determine whether it does anything useful. Rotations
alone would be useless---without the CX gates the circuit acts independently on
each qubit, making the output a deterministic function of each pixel separately,
with no interaction between pixels, which is the entire purpose of a convolution.
And the ring closure (last qubit back to qubit 0) allows information to reach
every qubit within one layer, where a plain chain would leave qubit 0 unable to
influence qubit $n-1$ until the second layer.

\subsection{Encoding}

Each pixel value $v \in [0,255]$ becomes a rotation angle $\theta = \pi v/255$
applied to its own qubit:
\begin{equation}
R_Y(\theta)|0\rangle = \cos(\theta/2)|0\rangle + \sin(\theta/2)|1\rangle
\end{equation}
A dark pixel remains near $|0\rangle$; a bright one rotates toward $|1\rangle$.
Because each qubit is rotated independently, the four-qubit input is a
\textbf{product state} with no entanglement. All entanglement in the output
originates in the filter's CX gates.

\section{Making Simulation Tractable}

A direct implementation submits one simulator job per patch. At 16 patches per
character and roughly 6{,}000 characters, that is approximately 96{,}000 jobs
per pass over the dataset---several hours, which makes iteration impossible.

Because the filter is \emph{fixed}, its unitary $U$ is identical for every patch.
It is extracted once and applied as a batched matrix multiplication:
\begin{equation}
|\psi_{\text{out}}\rangle = U\,|\psi_{\text{in}}\rangle
\end{equation}

\honestnote{%
\textbf{This is not an approximation.} It is the same linear algebra the
simulator performs, without per-job overhead and without shot noise. The claim
was verified rather than asserted: agreement with the shot-based implementation
was measured at a maximum deviation of \textbf{0.0028} across random patches at
200{,}000 shots, which is within sampling error. Processing time fell from hours
to \textbf{0.15 seconds for 2{,}000 characters}.}

\section{Readout Design}

The obvious readout is per-qubit $P(\text{measure }1)$: four values per patch.

But a 4-qubit circuit produces a distribution over 16 outcomes, and keeping only
four marginals discards every \emph{correlation} the entangling layer created---
precisely the component a classical product-state model could not have produced.
Adding pairwise correlations
\begin{equation}
\langle Z_i Z_j\rangle = \sum_k P(k)\, z_i(k)\, z_j(k)
\end{equation}
recovers six additional values per patch, taking the readout from 4 to 10 values
and the per-character feature vector from 64 to 160.

\begin{table}[h]
\centering
\caption{Effect of readout design, all tiers pooled, matched classifier.}
\begin{tabular}{@{}lrr@{}}
\toprule
Readout & Feature dimension & Accuracy \\
\midrule
Marginals only & 64 & 86.7\% \\
Marginals $+$ $\langle Z_i Z_j\rangle$ & 160 & \textbf{92.2\%} \\
\bottomrule
\end{tabular}
\end{table}

The 5.5-point gain from readout design alone exceeds any difference measured
between the quantum and classical filters. For quanvolutional approaches
generally, how much of the output distribution is retained appears to matter more
than which circuit produced it.

\section{Comparison Against Classical Baselines}

The comparison holds feature dimensionality fixed. Without that, a win would be
indistinguishable from simply granting one method a wider representation. The
classical control uses ten random filters over the same $2\times2$ patches at the
same stride, producing exactly 160 features.

\begin{table}[h]
\centering
\caption{Character accuracy, all tiers pooled.}
\begin{tabular}{@{}lrr@{}}
\toprule
Extractor & Dimension & Accuracy \\
\midrule
Quanvolutional, marginals & 64 & 86.7\% \\
Quanvolutional, $+\langle Z_iZ_j\rangle$ & 160 & 92.2\% \\
Classical convolution, dimension-matched & 160 & 91.9\% \\
Raw pixels, no feature extraction & 64 & \textbf{93.6\%} \\
\bottomrule
\end{tabular}
\end{table}

Repeating over 30 paired train--test splits establishes which differences are
real: the quantum-versus-classical gap is not distinguishable from noise
($p \approx 0.2$--$0.3$, the quantum layer winning approximately half the
splits), while the raw-pixel lead holds in all 30 ($p < 10^{-10}$).

\honestnote{%
\textbf{The quantum layer does not beat the classical baselines.} At matched
dimensionality it ties a random classical convolution, and both lose to raw
pixels at \emph{lower} dimensionality. The relevant comparison is therefore not
quantum versus classical but untrained random feature map versus no feature map:
at $8\times8$ the input is already close to information-minimal, so any fixed
random projection discards more than it contributes. That the classical control
fails identically is what identifies the cause as the untrained random character
of the filter rather than anything quantum.}

\section{Variational Circuit}

The brief specified a \emph{variational} circuit, so the filter's eight $R_Y$
angles were optimised against classification loss by gradient-free Powell search,
starting from the untrained baseline's exact angle vector so that any improvement
is attributable to optimisation alone. The classical control was trained
identically---training only the quantum side would have inverted the fairness the
untrained comparison was designed to preserve.

\begin{table}[h]
\centering
\caption{Variational training on a sealed test partition of 1{,}807 crops.}
\begin{tabular}{@{}lrrr@{}}
\toprule
Configuration & Trainable parameters & Accuracy & $\Delta$ \\
\midrule
Quantum, untrained & --- & 91.7\% & --- \\
Quantum, trained & 8 & 90.6\% & $-1.1$ pt \\
Classical, untrained & --- & 92.0\% & --- \\
Classical, trained & 40 & 92.1\% & $+0.1$ pt \\
Raw pixels & --- & \textbf{93.6\%} & --- \\
\bottomrule
\end{tabular}
\end{table}

Training does not close the gap, which strengthens rather than weakens the
preceding section: the deficit is not caused by the filters being untrained,
because training them does not help.

\subsection{A Correction Worth Recording}

An earlier version of this experiment reported a $+1.4$-point gain for the
trained quantum filter, bringing it level with raw pixels. That result was
wrong.

The angle search had been run over the full dataset, and the final evaluation
then used a fresh random split of that \emph{same} dataset---so the angles had
been selected using rows that later appeared in the test set. Re-running with a
sealed 30\% partition that the search never sees reduced the gain from $+1.4$
points to approximately zero.

The inflated figure is recorded here rather than discarded, because the
difference between the two is the difference between a false headline claim and
the actual finding.

\subsection{One Property That Survives}

The quantum filter produces its 160-dimensional feature space from
\textbf{8 parameters} against the classical control's \textbf{40}. At the
untrained baseline the two perform equivalently, so that $5\times$ parameter
economy is real and costs nothing there. It does not survive this training
protocol, and should not be reported without that qualification.

\section{Deliverable Status}

Complete. Quanvolutional layer implemented and verified against the reference
simulator; readout design measured; dimension-matched classical comparison
performed with significance testing; variational training implemented for both
quantum and classical filters.

\vspace{0.3cm}
\noindent\textbf{Artifacts:} \texttt{integration/features.py},
\texttt{track\_a\_vision/quanvolutional/train\_filter.py},
\texttt{notebooks/02\_quanvolutional\_features.ipynb},
\texttt{benchmarking/seed\_sweep.py}

\end{document}
